\documentclass{beamer}

\usetheme{Madrid}
\usepackage{amsmath,amssymb}
\usepackage{physics}
\usepackage{listings}
\usepackage{xcolor}

% ---------------------------
% Listings (Python)
% ---------------------------
\lstset{
  language=Python,
  basicstyle=\ttfamily\small,
  keywordstyle=\color{blue},
  commentstyle=\color{gray},
  stringstyle=\color{purple},
  showstringspaces=false,
  breaklines=true,
  frame=single,
  columns=fullflexible
}

\title{Entropy, Von Neumann Entropy, and Entanglement}
\author{Lecture Notes + Worked Numerical Examples}
\date{}

\begin{document}

%-------------------------------------------------
\begin{frame}
\titlepage
\end{frame}

%-------------------------------------------------
\begin{frame}{Outline}
\begin{enumerate}
    \item Shannon entropy (classical) and basic properties
    \item Density matrices and von Neumann entropy (quantum)
    \item Bell states and entanglement entropy
    \item Worked numerical examples (Python)
    \item Code slides: entropy computation and partial trace
\end{enumerate}
\end{frame}

%=================================================
\section{Classical entropy}

\begin{frame}{Entropy: Basic Idea}
\begin{itemize}
    \item Entropy quantifies \textbf{uncertainty} or \textbf{information content}.
    \item For a single event of probability $p$, the information content is
    \[
    I(p) = -\log_2 p.
    \]
    \item Rare events ($p\to 0$) carry large information; certain events ($p=1$) carry zero information.
\end{itemize}
\end{frame}

\begin{frame}{Shannon entropy}
For a discrete random variable $X$ with outcomes $x$:
\[
H(X) = -\sum_x p(x)\log_2 p(x).
\]

Binary entropy (outcomes with probabilities $p$ and $1-p$):
\[
H_2(p) = -p\log_2 p - (1-p)\log_2(1-p).
\]

\begin{itemize}
    \item $H_2(0)=H_2(1)=0$
    \item $H_2(1/2)=1$ bit (maximum uncertainty)
\end{itemize}
\end{frame}

%=================================================
\section{Quantum entropy}

\begin{frame}{Density matrices}
A quantum state can be:
\begin{itemize}
    \item \textbf{Pure:} $\rho = \ket{\psi}\bra{\psi}$, with $\rho^2=\rho$
    \item \textbf{Mixed:} $\rho=\sum_j p_j\ket{\psi_j}\bra{\psi_j}$, with $\rho^2\neq \rho$
\end{itemize}

Properties:
\[
\rho^\dagger=\rho,\quad \rho\succeq 0,\quad \Tr(\rho)=1.
\]
\end{frame}

\begin{frame}{Von Neumann entropy}
Quantum analogue of Shannon entropy:
\[
S(\rho) = -\Tr\!\left(\rho \log_2 \rho\right).
\]

If $\rho$ has eigenvalues $\{\lambda_i\}$:
\[
S(\rho)= -\sum_i \lambda_i \log_2 \lambda_i.
\]

\begin{itemize}
    \item Pure state $\Rightarrow \{\lambda\}=\{1,0,\dots\} \Rightarrow S(\rho)=0$
    \item Maximally mixed in $d$ dimensions: $\rho=\frac{I_d}{d} \Rightarrow S(\rho)=\log_2 d$
\end{itemize}
\end{frame}

%=================================================
\section{Bell states and entanglement entropy}

\begin{frame}{Bell states (two qubits)}
The four Bell states:
\begin{align*}
\ket{\Phi^+} &= \frac{1}{\sqrt{2}}(\ket{00}+\ket{11}), &
\ket{\Phi^-} &= \frac{1}{\sqrt{2}}(\ket{00}-\ket{11}),\\
\ket{\Psi^+} &= \frac{1}{\sqrt{2}}(\ket{01}+\ket{10}), &
\ket{\Psi^-} &= \frac{1}{\sqrt{2}}(\ket{01}-\ket{10}).
\end{align*}

They form an orthonormal basis of $\mathbb{C}^2\otimes \mathbb{C}^2$.
\end{frame}

\begin{frame}{Entanglement entropy for pure bipartite states}
For a pure state $\ket{\psi}_{AB}$ with $\rho_{AB}=\ket{\psi}\bra{\psi}$, define reduced states:
\[
\rho_A = \Tr_B(\rho_{AB}),\qquad \rho_B=\Tr_A(\rho_{AB}).
\]

\textbf{Entanglement entropy}:
\[
S_A := S(\rho_A),\qquad S_B := S(\rho_B).
\]

For pure bipartite states:
\[
S(\rho_A)=S(\rho_B).
\]
Moreover, $S(\rho_A)=0$ iff $\ket{\psi}$ is separable.
\end{frame}

\begin{frame}{Worked derivation: Bell state reduced density matrix}
Take $\ket{\Phi^+}=\frac{1}{\sqrt{2}}(\ket{00}+\ket{11})$.
Then
\[
\rho_{AB}=\ket{\Phi^+}\bra{\Phi^+}
=\frac{1}{2}\Big(\ket{00}\bra{00}+\ket{00}\bra{11}+\ket{11}\bra{00}+\ket{11}\bra{11}\Big).
\]

Partial trace over $B$ using $\Tr_B(\ket{ab}\bra{cd})=\braket{d}{b}\ket{a}\bra{c}$:
\[
\rho_A=\Tr_B(\rho_{AB})=\frac{1}{2}\Big(\ket{0}\bra{0}+\ket{1}\bra{1}\Big)=\frac{I_2}{2}.
\]

Eigenvalues of $\rho_A$ are $\{1/2,1/2\}$, so
\[
S(\rho_A)= -2\cdot\frac{1}{2}\log_2\left(\frac{1}{2}\right)=1~\text{bit}.
\]
Thus Bell states are \textbf{maximally entangled}.
\end{frame}

%=================================================
\section{Worked numerical examples (Python)}

\begin{frame}[fragile]{Numerical example 1: binary Shannon entropy curve}
Compute $H_2(p)$ for representative $p$ values.
\begin{lstlisting}
import math

def H2(p):
    eps = 1e-15
    p = min(max(p, eps), 1.0-eps)
    return -p*math.log(p, 2) - (1-p)*math.log(1-p, 2)

for p in [0.0, 0.1, 0.25, 0.5, 0.75, 0.9, 1.0]:
    # clamp for endpoints
    pp = min(max(p, 1e-15), 1-1e-15)
    print(p, H2(pp))
\end{lstlisting}

\textbf{Check:} maximum near $p=0.5$, and near-zero for $p\approx 0$ or $1$.
\end{frame}

\begin{frame}[fragile]{Numerical example 2: von Neumann entropy from eigenvalues (2x2)}
For a $2\times 2$ density matrix, eigenvalues can be computed analytically.
\begin{lstlisting}
import math, cmath

def eigvals_2x2(rho):
    a, b = rho[0]
    c, d = rho[1]
    tr = a + d
    det = a*d - b*c
    disc = tr*tr - 4*det
    root = cmath.sqrt(disc)
    lam1 = 0.5*(tr + root)
    lam2 = 0.5*(tr - root)
    # return real parts (should be real for Hermitian rho)
    return [lam1.real, lam2.real]

def S_vn_from_eigs(lams):
    S = 0.0
    for lam in lams:
        lam = max(lam, 0.0)
        if lam > 1e-15:
            S -= lam*math.log(lam, 2)
    return S

rhoA = [[0.5, 0.0],
        [0.0, 0.5]]
print("eigvals:", eigvals_2x2(rhoA))
print("S(rhoA):", S_vn_from_eigs(eigvals_2x2(rhoA)))
\end{lstlisting}

For Bell-state reductions, you should get $S(\rho_A)=1$.
\end{frame}

\begin{frame}[fragile]{Code slide: build Bell states and compute entanglement entropy}
We compute $\rho_A=\Tr_B(\rho_{AB})$ for Bell states and evaluate $S(\rho_A)$.
\begin{lstlisting}
import math

def outer(v):
    n = len(v)
    return [[v[i]*v[j].conjugate() for j in range(n)] for i in range(n)]

def partial_trace_B(rho):
    # 2 qubits, basis: |00>,|01>,|10>,|11>
    rhoA = [[0j, 0j],[0j, 0j]]
    for a in (0,1):
        for c in (0,1):
            s = 0j
            for b in (0,1):
                i = 2*a + b
                j = 2*c + b
                s += rho[i][j]
            rhoA[a][c] = s
    return rhoA

# Bell states as length-4 vectors
inv = 1/math.sqrt(2)
Phi_plus  = [inv, 0j, 0j, inv]
Phi_minus = [inv, 0j, 0j, -inv]
Psi_plus  = [0j, inv, inv, 0j]
Psi_minus = [0j, inv, -inv, 0j]

for name, psi in [("Phi+", Phi_plus), ("Phi-", Phi_minus),
                  ("Psi+", Psi_plus), ("Psi-", Psi_minus)]:
    rhoAB = outer(psi)
    rhoA  = partial_trace_B(rhoAB)
    # convert to plain real for the eig solver
    rhoA_real = [[rhoA[0][0].real, rhoA[0][1].real],
                 [rhoA[1][0].real, rhoA[1][1].real]]
    lams = eigvals_2x2(rhoA_real)
    print(name, "rhoA eigenvalues:", lams, "S(rhoA)=", S_vn_from_eigs(lams))
\end{lstlisting}
Expected output: eigenvalues $(1/2,1/2)$ and entropy $1$ for all Bell states.
\end{frame}

\begin{frame}[fragile]{Numerical example 3: partially entangled state $\cos\theta|00\rangle+\sin\theta|11\rangle$}
This illustrates how entanglement entropy varies continuously from 0 to 1.
\begin{lstlisting}
import math

def psi_theta(theta):
    return [math.cos(theta), 0j, 0j, math.sin(theta)]

print("theta   S(rhoA)")
for k in range(11):
    theta = (math.pi/2) * k/10
    psi = psi_theta(theta)
    rhoAB = outer(psi)
    rhoA = partial_trace_B(rhoAB)
    rhoA_real = [[rhoA[0][0].real, rhoA[0][1].real],
                 [rhoA[1][0].real, rhoA[1][1].real]]
    S = S_vn_from_eigs(eigvals_2x2(rhoA_real))
    print(f"{theta:6.3f}  {S:8.6f}")
\end{lstlisting}
Check: $S=0$ at $\theta=0$ and $\theta=\pi/2$, and $S=1$ at $\theta=\pi/4$.
\end{frame}

%=================================================
\section{Optional: with NumPy (compact + scalable)}

\begin{frame}[fragile]{Optional compact implementation using NumPy}
If NumPy is allowed, computation becomes very compact:
\begin{lstlisting}
import numpy as np

def vn_entropy(rho):
    # eigenvalues of Hermitian matrix
    w = np.linalg.eigvalsh(rho)
    w = np.maximum(w, 0.0)
    w = w[w > 1e-15]
    return -np.sum(w*np.log2(w))

# Bell state |Phi+>
psi = np.array([1,0,0,1], dtype=complex)/np.sqrt(2)
rhoAB = np.outer(psi, psi.conj())

# partial trace over B
rhoA = np.zeros((2,2), dtype=complex)
for a in [0,1]:
    for c in [0,1]:
        for b in [0,1]:
            rhoA[a,c] += rhoAB[2*a+b, 2*c+b]

print(vn_entropy(rhoA))  # should print 1.0
\end{lstlisting}
\end{frame}

%=================================================
\section{Summary}

\begin{frame}{Summary}
\begin{itemize}
    \item Shannon entropy $H(X)$ measures classical uncertainty.
    \item Von Neumann entropy $S(\rho)$ measures quantum mixedness.
    \item For a pure bipartite state, $S(\rho_A)$ quantifies \textbf{entanglement}.
    \item Bell states have reduced density matrix $\rho_A=I/2$ and entanglement entropy $S=1$ bit.
    \item Python examples confirm theory and give templates for larger simulations.
\end{itemize}
\end{frame}

\end{document}
